\documentclass[10pt,a4paper]{amsart}
\usepackage[margin=19mm]{geometry}
\usepackage{amsmath,amssymb,mathtools}
\usepackage[scr=rsfs,cal=euler]{mathalfa}
\usepackage{xfrac}
\usepackage[unicode,hidelinks,bookmarks=false]{hyperref}
\newcommand{\ie}{i.e.\ }
\newcommand{\cf}{cf.\ }
\newcommand{\resp}{respectively\ }
\newcommand{\Subsection}{\subsection}
\providecommand{\nobreakdash}{\nobreak\hbox{-}}
\setlength{\emergencystretch}{2em}
\multlinegap0pt
\usepackage{mathtools}
\usepackage{enumitem}
\newtheorem{theorem}{Theorem}
\newtheorem{lemma}[theorem]{Lemma}
\newtheorem{proposition}[theorem]{Proposition}
\newtheorem{corollary}[theorem]{Corollary}
\newtheorem{definition}[theorem]{Definition}
\newtheorem{remark}[theorem]{Remark}
\newtheorem{conjecture}[theorem]{Conjecture}
\newcommand{\Fpm}{\{-1,+1\}}
\newcommand{\jc}{j}
\newcommand{\Jc}{J}
\begin{document}
\title[A divisibility theorem for odd J-characteristics of two-leve]{A divisibility theorem for odd $J$-characteristics of two-level designs}
\author{Pieter Thijs Eendebak}
\date{}
\maketitle
\noindent\emph{Source and licence.} A divisibility theorem for odd $J$-characteristics of two-level designs. Version arXiv:2607.09478v1 (CC BY 4.0). This material is distributed under the Creative Commons Attribution 4.0 International licence, http://creativecommons.org/licenses/by/4.0/, as recorded in the arXiv OAI licence metadata for this exact version and shown on the abstract page https://arxiv.org/abs/2607.09478v1. Full rights notice: licenses/arxiv-2607.09478-CC-BY-4.0.txt.\par\smallskip
\noindent\textit{Editorial scope.} Counted unit: the original \texttt{theorem} environment \texttt{thm:divis} at source lines 144--150 together with its complete proof at lines 152--170; the delivered document keeps source lines 109--171 so that every referenced label and every citation resolves inside it. Native editable source is the paper's own concatenated TeX; the AMS wrapper is editorial formatting only. No publisher class, upstream script or unrelated artwork is redistributed.
\begin{lemma}[Orthogonality]\label{lem:orth}
For $x,y\in\Fpm^n$,
\begin{equation}\label{eq:orth}
  \sum_{s\subseteq[n]} \chi_s(x)\,\chi_s(y)
  =\begin{cases} 2^n & x=y,\\ 0 & x\ne y.\end{cases}
\end{equation}
\end{lemma}

\begin{proof}
$\chi_s(x)\chi_s(y)=\prod_{i\in s} x_i y_i$, so
$\sum_{s\subseteq[n]}\chi_s(x)\chi_s(y)=\sum_{s\subseteq[n]}\prod_{i\in s}x_iy_i=\prod_{i=1}^n(1+x_iy_i)$,
the last step by expanding the product over all subsets. Each factor is $2$ when
$x_i=y_i$ and $0$ otherwise.
\end{proof}

\begin{lemma}[Inversion]\label{lem:inv}
For every $f:\Fpm^n\to\mathbb Z$ and every $x\in\Fpm^n$,
\begin{equation}\label{eq:inv}
  2^n\, f(x)=\sum_{s\subseteq[n]} \jc(s)\,\chi_s(x).
\end{equation}
\end{lemma}

\begin{proof}
Expand $\jc(s)=\sum_y f(y)\chi_s(y)$, swap sums, and apply
Lemma~\ref{lem:orth}:
\[
\sum_s\jc(s)\chi_s(x)
=\sum_y f(y)\sum_s\chi_s(y)\chi_s(x)
=\sum_y f(y)\cdot 2^n[y{=}x]
=2^n f(x).\qedhere
\]
\end{proof}

\section{The divisibility theorem}


\begin{theorem}[Divisibility of the top odd characteristic]\label{thm:divis}
Let $n\ge 1$ be \emph{odd} and let $f:\Fpm^n\to\mathbb Z$. If
\begin{equation}\label{eq:divhyp}
  \jc(s)=0 \quad\text{for every }s\subseteq[n]\text{ with $|s|$ odd and }s\ne[n],
\end{equation}
then $2^{\,n-1}\mid \jc([n])$.
\end{theorem}

\begin{proof}
Apply Lemma~\ref{lem:inv} at $x$ and at $-x$. Since
$\chi_s(-x)=(-1)^{|s|}\chi_s(x)$, subtracting gives
\[
  2^n\bigl(f(x)-f(-x)\bigr)
   = \sum_{s\subseteq[n]} \jc(s)\bigl(1-(-1)^{|s|}\bigr)\chi_s(x)
   = 2\!\!\sum_{|s|\text{ odd}}\!\!\jc(s)\,\chi_s(x),
\]
because $1-(-1)^{|s|}$ is $0$ for even $|s|$ and $2$ for odd $|s|$. By
\eqref{eq:divhyp} the only odd-cardinality $s$ with $\jc(s)\ne 0$ is
$s=[n]$ --- here we use that $n$ is odd, so $[n]$ itself has odd cardinality.
The right-hand sum collapses to $\jc([n])\chi_{[n]}(x)$, and dividing by~$2$:
\begin{equation}\label{eq:key}
  2^{\,n-1}\bigl(f(x)-f(-x)\bigr)=\jc([n])\,\chi_{[n]}(x).
\end{equation}
Take $x=(+1,\dots,+1)$, so $\chi_{[n]}(x)=1$; the left-hand side is
$2^{\,n-1}$ times the integer $f(x)-f(-x)$, hence
$2^{\,n-1}\mid \jc([n])$.
\end{proof}


\end{document}
